% Independent derivation, 2026-09-24. FSP1--20.
\section{The ungraded presence unit, spherical algebra, and arithmetic realization}
\label{sec:f1-spherical-parity-base-change}

\paragraph{Original sources and exact scope.}
Alain Connes and Caterina Consani,
\href{https://arxiv.org/abs/2606.06604v1}{\emph{On the Absolute
Geometry of $\operatorname{Spec}\mathbb Z$ and the Fargues--Fontaine
curve}}, original \nolinkurl{FF.tex}, line 145, identify their
$\mathbb F_1$ with the spherical algebra of the pointed
multiplicative monoid $\{0,1\}$. Lines 343--371 describe its
arithmetic stalks; lines 374--404 describe morphisms through
the underlying multiplicative monoid, with the exponent-zero
condition $\chi(0)=1$. The source SHA-256 is
\begin{center}\small
\nolinkurl{14c17a95040be7bfc57bb5f87e58907a05cacd0bcb1b6967b2892cf3d9bc1229}.
\end{center}
Their earlier
\href{https://arxiv.org/abs/1502.05585v2}{\emph{Absolute algebra
and Segal's $\Gamma$-rings}}, original \nolinkurl{F1_corr.tex},
gives the needed definitions and comparisons: lines 416--438
define the category and the sum-over-fibres functor; lines
464--501 define the sphere and spherical monoid algebra;
Proposition \texttt{monomono}, lines 504--510, proves the
monoid adjunction; Lemma \texttt{sssalg}, lines 519--532,
defines the semiring construction and its unit. Lines 544--560
give the two projections and folding definition of addition.
The Boolean construction and comparison with $\mathbb F_2$
are at lines 568--580, and the absence of two simultaneous
nonzero inputs in a spherical monoid algebra is at lines
865--880. Its source SHA-256 is
\begin{center}\small
\nolinkurl{3b1319ec1df4ed6e2412ace960c2813b683a36c88509234c9abd0a39f98e8913}.
\end{center}
These precise passages were read in original author TeX for
this derivation. The calculations below instantiate their
existing constructions; the source mathematics is not claimed
as a new discovery. The support operations remain those of
The Clankers,
\href{https://github.com/KokunoYumeto/zeta-function-research-reader/blob/a2851f9eb352e7e4376bc629de86fa4ed9c1c434/workbenches/splitzero-tandem/foundations/split-support-geometry-v11/split_support_geometry_arithmetic_curve_v11.tex#L1226}
{\emph{Split Support Geometry}, Definition \texttt{def:lattice-split}}.

\subsection{Presence has a multiplicative unit before arithmetic parity}

Keep the user's presence unit $\tau$ and denote the external
absent point by $0_M$. Define the pointed monoid
\begin{equation}\tag{FSP1}
M=\{0_M,\tau\},\qquad
\tau^2=\tau,\qquad
0_M\tau=\tau0_M=0_M^2=0_M.
\end{equation}
Here $\tau$ is the multiplicative unit and $0_M$
the absorber. No addition and no parity function are part of
this multiplicative structure. The named map
$j:M\to\{0,1\}_{\rm CC}$ with $j(0_M)=0$ and $j(\tau)=1$
is an isomorphism of pointed multiplicative monoids: its four
products agree with the four products of FSP1, and the map is
bijective. Its target is Connes--Consani's pointed monoid.
The singleton $\{\tau\}$ supplies its
unit monoid; $M$ is that singleton with an external basepoint
adjoined. This is the exact structure used here, not an
identification of $\tau$ with a pre-existing integer.

Let $\Gamma^{\mathrm{op}}$ have objects
$k_+=\{0,1,\ldots,k\}$ and all basepoint-preserving maps,
as in the original source. The spherical monoid algebra is
\begin{equation}\tag{FSP2}
\mathbb S M(k_+)=M\wedge k_+
=\{0,\tau_1,\ldots,\tau_k\},\qquad
\mathbb S M(f)(\tau_j)=
\begin{cases}\tau_{f(j)}&f(j)\ne0,\\0&f(j)=0.\end{cases}
\end{equation}
Here $\tau_j$ is the class of $(\tau,j)$; all pairs
with either coordinate a basepoint are identified with $0$.
Composition follows immediately from composition of $f$.
The product sends $(\tau_i,\tau_j)$ to $\tau_{(i,j)}$ at
$k_+\wedge l_+$, with every product containing a basepoint
sent to the basepoint. Associativity is associativity of the
label pair $(i,j)$, and the one available unit is
$\tau$. The maps $\tau_j\mapsto j$ and $0\mapsto0$
give an isomorphism of spherical algebras
\begin{equation}\tag{FSP3}
\mathbb S M\cong\mathbb S,\qquad\mathbb S(k_+)=k_+.
\end{equation}
Thus this is the specified Connes--Consani spherical
$\mathbb F_1$ model. Its description requires no parity assignment.

\subsection{Every coefficient realization, with all coordinates}

For any commutative unital semiring $A$, let
\begin{equation}\tag{FSP4}
HA(k_+)=A^k,\qquad
(HA(f)(a))_r=\sum_{j\in f^{-1}(r)}a_j\quad(1\le r\le l).
\end{equation}
The empty sum is $0_A$; the basepoint is the all-$0_A$
tuple. Associativity and commutativity of addition prove
functoriality: the sum for a composite groups exactly the
same finite terms according to the intermediate fibre.
Multiplication on two inputs is the tensor-coordinate rule
$(a,b)\mapsto(a_i b_j)_{(i,j)}$ at $k_+\wedge l_+$.
Distributivity proves its naturality, and semiring
associativity and the unit prove the spherical algebra laws.

The exact unit realization is
\begin{equation}\tag{FSP5}
\eta_{A,k}:\mathbb S M(k_+)\longrightarrow HA(k_+),\qquad
0\longmapsto(0_A,\ldots,0_A),\qquad
\tau_j\longmapsto\delta_j^A,
\end{equation}
where $\delta_j^A$ has entry $1_A$ at $j$ and $0_A$ at
all other coordinates. Under (FSP4), a delta vector stays
a delta at $f(j)$ if that index is nonzero, and becomes
the zero tuple otherwise. Hence (FSP5) is natural for
every pointed map. A product of two delta vectors is the
delta at $(i,j)$, proving multiplicativity and preservation
of the spherical unit. This supplies a complete spherical
algebra morphism, not merely a map at $1_+$.

Specializing $A=\mathbb Z$ gives
\begin{equation}\tag{FSP6}
\eta_{\mathbb Z,k}(\tau_j)=e_j\in\mathbb Z^k.
\end{equation}
The image of the ungraded presence unit at level one is
integer $1$. The map and codomain give that image its
arithmetic meaning. They do not retroactively add an
integer grading to the definition (FSP1).

\subsection{What the fold does, and what its domain contains}

Let $p_1,p_2:2_+\to1_+$ select the first or the second
nonzero coordinate and send the other to the basepoint.
Let $\nabla:2_+\to1_+$ send both $1$ and $2$ to $1$.
Then
\begin{equation}\tag{FSP7}
\begin{array}{c|ccc}
z\in\mathbb S M(2_+)&0&\tau_1&\tau_2\\\hline
\mathbb S M(p_1)z&0&\tau&0\\
\mathbb S M(p_2)z&0&0&\tau\\
\mathbb S M(\nabla)z&0&\tau&\tau
\end{array}
\end{equation}
There is no $z$ whose two projections are both
$\tau$. In the source's generalized-addition
definition, this says the set of witnesses for
$\tau\oplus\tau$ is empty.
It does not say the fold maps a simultaneous pair to
$\tau$: no such pair lies in its domain.

In contrast $H\mathbb Z(2_+)=\mathbb Z^2$ contains the
additional tuple $(1,1)$, and
\begin{equation}\tag{FSP8}
H\mathbb Z(\nabla)(1,1)=1+1=2.
\end{equation}
The image of $\eta_{\mathbb Z,2}$ is precisely
$\{(0,0),(1,0),(0,1)\}$ and does not contain $(1,1)$.
Thus (FSP7) and (FSP8) are compatible under an exact
morphism. They describe different available inputs,
with all coordinates displayed.

\subsection{Boolean addition is an explicit extension of the source}

If one adds the law
$\tau+\tau=\tau$ and makes
$0_M$ an additive identity, the result is the Boolean
semiring $B_M\cong\mathbf B$, not merely (FSP1).
Its spherical algebra is
\begin{equation}\tag{FSP9}
HB_M(k_+)=\{0_M,\tau\}^k,\qquad
HB_M(\nabla)(\tau,\tau)
   =\tau.
\end{equation}
It receives the same delta construction
$\eta_{B_M}:\mathbb S M\to HB_M$.
There is no unital spherical algebra morphism
$HB_M\to H\mathbb Z$. To prove this directly, its
level-one map would send $0_M$ to $0$ and
$\tau$ to $1$. Naturality under $p_1,p_2$
then forces $(\tau,\tau)$ to map
to $(1,1)$. Naturality under $\nabla$ would give
$1=2$, contradicting integer arithmetic.
The same argument forbids a unital map $HB_M\to H\mathbb F_2$:
the fold there gives $1+1=0$, not $1$.
The actual comparison is the pair of proved maps
\begin{equation}\tag{FSP10}
HB_M\ \longleftarrow\ \mathbb S M\ \longrightarrow\ H\mathbb Z.
\end{equation}
Their existence retains the common multiplicative source
without falsely identifying its two additive extensions.

There is also an exact action, on the entire integer
receiver, of the pointed multiplicative monoid underlying
$B_M$:
\begin{equation}\tag{FSP11}
\rho:B_M^{\mathrm{mult}}\longrightarrow
\operatorname{End}_{+}(\mathbb Z)^{\mathrm{mult}},\qquad
\rho(0_M)=0,\qquad\rho(\tau)=\operatorname{id}.
\end{equation}
Here the multiplication in the endomorphism monoid is
composition; each of its displayed operators preserves
receiver addition. Its pointed multiplicative table is
exactly (FSP1), proving the action. It does not preserve
addition of scalars, and makes the two calculations explicit:
\begin{equation}\tag{FSP12}
(\tau+\tau)\mathbin{\cdot}1
 =\tau\mathbin{\cdot}1=1,
\qquad
\tau\mathbin{\cdot}1+
\tau\mathbin{\cdot}1=1+1=2.
\end{equation}
The first equality uses Boolean substitution followed
by the action. The second is addition in the integer
receiver. No distributive equality between these expressions
is asserted or needed. If scalar distributivity were imposed,
it would force $x=x+x$ for every receiver element; on an
abelian-group receiver that condition forces $x=0$.
The existing nontrivial monoid action (FSP11) remains valid.

\subsection{The full support carrier and its distinct zero}

For any bounded distributive lattice $L$, retain
\begin{equation}\tag{FSP13}
\begin{gathered}
A_L=G_L(\mathbb Z)=\{(0,\lambda):\lambda\in L\}
           \cup\{(n,1_L):n\in\mathbb Z\},\\
(n,\lambda)+(m,\mu)=(n+m,\lambda\vee\mu),\qquad
(n,\lambda)(m,\mu)=(nm,\lambda\wedge\mu),\\
0_{A_L}=\tau_L=(0,0_L),\quad e_L=(0,1_L),\quad
1_{A_L}=(1,1_L).
\end{gathered}
\end{equation}
The exact realization (FSP5) in these coefficients has
\begin{equation}\tag{FSP14}
\eta_{A_L,k}(\tau_j)_i=
\begin{cases}(1,1_L)&i=j,\\\tau_L&i\ne j,\end{cases}
\qquad \eta_{A_L,k}(0)=(\tau_L,\ldots,\tau_L).
\end{equation}
Thus the presence unit $\tau$ is mapped to
the ambient multiplicative unit, and the absent point
$0_M$ to the ambient additive zero. The original absorber
$\tau_L$ is not renamed a multiplicative unit. Mapping
$\tau$ to $\tau_L$ instead would fail unit
preservation, because $(1,1_L)\ne(0,0_L)$.
There is nevertheless an exact corner realization using
the old $\tau_L$ itself as a unit. Adjoin a strict new
bottom $*$ to $L$, put $A_L^+=G_{L^+}(\mathbb Z)$,
and write $z_*=(0,*)$. Then
\begin{equation}\tag{FSP14a}
C=\tau_L A_L^+=\{z_*,\tau_L\},\qquad
M\longrightarrow C^{\mathrm{mult}},\quad
0_M\longmapsto z_*,\quad\tau\longmapsto\tau_L.
\end{equation}
This is an isomorphism of pointed multiplicative monoids:
$\tau_L^2=\tau_L$, $\tau_Lz_*=z_*$, and $z_*^2=z_*$.
The corner also has Boolean addition with zero $z_*$
and unit $\tau_L$. It therefore gives the exact map
$\mathbb S M\to HC$ by (FSP5).
The inclusion $C\to A_L^+$ preserves addition,
multiplication, and zero, but not the multiplicative unit;
its induced $HC\to HA_L^+$ is correspondingly a
multiplicative natural transformation that fails unit
preservation as a spherical algebra map. In the other
direction, $x\mapsto\tau_Lx$ is a unital semiring map
$A_L^+\to C$: every old element maps to $\tau_L$,
and $z_*$ maps to $z_*$. It is precisely localization
at $\tau_L$, since an invertible idempotent must be
the unit and then $\tau_Lx=\tau_L$ forces every old
$x$ to that unit. The two elements remain distinct
under the displayed map to $C$.
Thus the old symbol can realize the presence unit in
its named corner or localization, with the complete
maps and their different unit conditions specified.
The complete functor on a tuple retains cancellations as
\begin{equation}\tag{FSP15}
\left(H A_L(f)((n_j,\lambda_j)_j)\right)_r
 =\left(\sum_{j\in f^{-1}(r)}n_j,
              \bigvee_{j\in f^{-1}(r)}\lambda_j\right),
\end{equation}
with empty sum $0$ and empty join $0_L$. In particular
the fold of $((1,1_L),(-1,1_L))$ is $e_L$, retaining
top support, and the fold of $((1,1_L),(1,1_L))$ is
$(2,1_L)$.

Amplitude, support, and parity define exact coefficient
homomorphisms
\begin{equation}\tag{FSP16}
\begin{gathered}
p_L:A_L\to\mathbb Z,\quad(n,\lambda)\mapsto n,\qquad
\chi_L:A_L\to L,\quad(n,\lambda)\mapsto\lambda,\\
\pi_{2,L}:A_L\to G_L(\mathbb F_2),\quad
(n,\lambda)\mapsto(\overline n,\lambda),
\qquad p_2\pi_{2,L}=\pi_2p_L.
\end{gathered}
\end{equation}
All preserve both operations, zero, and unit by (FSP13).
Applying $H$ coordinatewise proves naturality using
(FSP15) and gives, in particular,
$Hp_L\eta_{A_L}=\eta_{\mathbb Z}$ and
$H\pi_{2,L}\eta_{A_L}=\eta_{G_L(\mathbb F_2)}$.
The full parity map still distinguishes every zero label;
supported even integers map to $e_L$.
If $L$ is nontrivial, there is no unital spherical algebra
map $H\mathbb Z\to HA_L$ serving as a reverse arithmetic
inclusion. Indeed naturality under coordinate projections
and folding forces its level-one map to be additive, and
unit preservation forces $1\mapsto(1,1_L)$. The image
of $-1$ would then have to satisfy
$(1,1_L)+x=\tau_L$, impossible because the support of
that sum is $1_L$. The forward amplitude map and the
spherical unit map remain the exact comparison.

The action (FSP11) also exists on the entire additive
monoid of $A_L$: the two operators are the constant
$\tau_L$ map and the identity. Each preserves addition,
their composition table is (FSP1), and all labels are
retained by the identity. Pointwise addition of two identity
operators sends $(n,\lambda)$ to $(2n,\lambda)$; this
equals the identity exactly on the complete zero fibre.
Thus the failure of scalar additivity and its fixed locus
are explicitly computed in the original coefficients.

\subsection{Free arithmetic base change and the origin of parity}

Form the ordinary monoid ring $\mathbb Z[M]$ with basis
$[0_M],[\tau]$ and multiplication $[a][b]=[ab]$.
Its multiplicative identity is $[\tau]$.
Contract the absorbing basis element by its generated ideal:
\begin{equation}\tag{FSP17}
\mathbb Z_0[M]=\mathbb Z[M]/([0_M])\cong\mathbb Z,
\qquad n[\tau]\longmapsto n.
\end{equation}
The ideal $([0_M])$ is exactly $\mathbb Z[0_M]$, since
multiplying either basis vector by $[0_M]$ gives $[0_M]$.
The quotient is therefore free of rank one with basis
$[\tau]$, and that basis vector is its unit;
this proves both injectivity and surjectivity of the
displayed ring map. Equivalently a pointed multiplicative
map $M\to R^{\mathrm{mult}}$ to any unital ring must send
$0_M$ to $0_R$ and $\tau$ to $1_R$; it extends
uniquely to the ring homomorphism $\mathbb Z\to R$.
This proves the free arithmetic base-change property directly.

For this realization the composite is
\begin{equation}\tag{FSP18}
M\longrightarrow\mathbb Z_0[M]\cong\mathbb Z
\longrightarrow\mathbb F_2,
\qquad 0_M\longmapsto0\longmapsto0,
\quad\tau\longmapsto1\longmapsto1.
\end{equation}
Thus the realized absent point has even integer image and
the realized unit has odd integer image. These are image
properties under named maps. The source (FSP1) has no
intrinsic integer parity predicate or additive cancellation
rule. Nor is (FSP17) a base change of the extra Boolean
addition: any additive unit-preserving map from $B_M$ to
a ring would impose $1+1=1$, forcing the zero ring.
Forgetting the extra Boolean addition before (FSP17)
is therefore an explicit change of category, not an
unmentioned equality of the two extensions.

\subsection{The map on Connes--Consani's actual stalk generators}

At a closed arithmetic prime $p$, write
$H_p^+=\mathbb Z[1/p]_{\ge0}$ and form the pointed
monoid
\begin{equation}\tag{FSP19}
N_p=\{0_{N_p}\}\sqcup\{T^a:a\in H_p^+\},\qquad
T^aT^b=T^{a+b},\qquad T^0=\tau.
\end{equation}
The source stalk is $\mathbb S N_p$, agreeing with the
spherical presentation in \nolinkurl{FF.tex}, Proposition
\texttt{prop12}. Every pointed multiplicative map
$c:N_p\to A_L$ with $c(T^0)=1_{A_L}$ determines the
exact spherical map
\begin{equation}\tag{FSP20}
\mathbb S N_p(k_+)\longrightarrow HA_L(k_+),\qquad
(T^a,j)\longmapsto
(\tau_L,\ldots,\tau_L,c(T^a),\tau_L,\ldots,\tau_L),
\quad 0\longmapsto(\tau_L,\ldots,\tau_L).
\end{equation}
Naturality follows because each input has only one active
coordinate, and products are preserved precisely by
$c(T^{a+b})=c(T^a)c(T^b)$. Conversely a spherical
algebra map restricts at $1_+$ to a pointed multiplicative
map. Every $(T^a,j)$ is the image of $T^a$ under the
pointed injection $1_+\to k_+$ selecting $j$, so
naturality forces exactly (FSP20); there is no additional
choice. This proves the monoid adjunction here explicitly.
Amplitude, support, and parity of every such stalk morphism
are then obtained by composing $c$ with the exact maps
(FSP16), at every coordinate. No assigned source parity
is needed, and no old zero label is discarded except
by the explicitly named amplitude projection.
